\documentclass[12pt]{article}
\usepackage[T1]{fontenc}
\usepackage[utf8]{inputenc}
\usepackage{lmodern}
\usepackage{textcomp}
\usepackage{enumitem}
\usepackage{geometry}
\geometry{margin=1in}
\begin{document}
\begin{center}
Answer Key - Physics - Undergraduate Level Assessment 4 \
\small (Answers are ordered exactly as in the question paper.)
\end{center}

\section*{Multiple Choice}
\begin{enumerate}[label=\arabic*.]
\item \textbf{Answer:} A
\\ \textit{Options:} A) gamma rays; B) microwaves; C) ultraviolet radiation; D) visible light
\\ \textit{Note:} Electromagnetic radiation emitted from a nucleus is most likely to be in the form of gamma rays because gamma rays are high-energy photons released during nuclear decay processes, reflecting the transition of the nucleus to a lower energy state.
\item \textbf{Answer:} D
\\ \textit{Options:} A) 2; B) 4; C) 6; D) 10
\\ \textit{Note:} The atom has filled n = 1 and n = 2 energy levels, which can hold a maximum of 2 electrons in n = 1 and 8 electrons in n = 2, resulting in a total of 10 electrons.
\item \textbf{Answer:} D
\\ \textit{Options:} A) 0.25 W; B) 0.5 W; C) 1 W; D) 4 W
\\ \textit{Note:} The rate of energy dissipation in a resistor is given by the formula P = V²/R, so if the voltage (V) is doubled, the power dissipated increases by a factor of four, resulting in a new rate of 4 W.
\item \textbf{Answer:} D
\\ \textit{Options:} A) Bosons have symmetric wave functions and obey the Pauli exclusion principle.; B) Bosons have antisymmetric wave functions and do not obey the Pauli exclusion principle.; C) Fermions have symmetric wave functions and obey the Pauli exclusion principle.; D) Fermions have antisymmetric wave functions and obey the Pauli exclusion principle.
\\ \textit{Note:} This statement is correct because fermions, which are particles with half-integer spin, require their wave functions to be antisymmetric under particle exchange, leading to the Pauli exclusion principle that prohibits identical fermions from occupying the same quantum state.
\item \textbf{Answer:} D
\\ \textit{Options:} A) 0.1 GeV/$c^2$; B) 0.2 GeV/$c^2$; C) 0.5 GeV/$c^2$; D) 1.0 GeV/$c^2$
\\ \textit{Note:} The rest mass can be determined using the energy-momentum relation \( E^2 = (pc)^2 + (m_0 c^2)^2 \), where substituting \( E = 5.0 \, \text{GeV} \) and \( p = 4.9 \, \text{GeV/c} \) leads to a calculation revealing that the rest mass \( m_0 \) is approximately \( 1.0 \, \text{GeV/c}^2 \).
\end{enumerate}

\section*{True / False}
\begin{enumerate}[label=\arabic*.]
\setcounter{enumi}{5}
\item \textbf{Answer:} False
\\ \textit{Options:} A) True; B) False
\\ \textit{Note:} A helium-neon laser, while useful for certain applications, is not the most versatile for spectroscopy across visible wavelengths because it has limited output wavelengths compared to other lasers, such as dye lasers or solid-state lasers, which can be tuned across a broader range of the visible spectrum.
\item \textbf{Answer:} False
\\ \textit{Options:} A) True; B) False
\\ \textit{Note:} The negative muon (mu^-) is a lepton, which are elementary particles that do not experience strong interactions, whereas mesons are composite particles made of quark-antiquark pairs that do participate in strong interactions, making their properties fundamentally different.
\item \textbf{Answer:} False
\\ \textit{Options:} A) True; B) False
\\ \textit{Note:} The angular magnification of a refracting telescope is given by the ratio of the focal lengths of the objective lens to the eyepiece lens, and with a typical arrangement, a 20 cm focal length eyepiece does not yield an angular magnification of 6, indicating the statement is false.
\item \textbf{Answer:} False
\\ \textit{Options:} A) True; B) False
\\ \textit{Note:} Phosphorus is not a suitable dopant for germanium to form an n-type semiconductor because phosphorus has a larger atomic size and different electronegativity compared to germanium, which can lead to poor solubility and reduced effectiveness in creating free electrons.
\item \textbf{Answer:} True
\\ \textit{Options:} A) True; B) False
\\ \textit{Note:} Gas lasers operate by utilizing transitions between the energy levels of free atoms in a gaseous medium because these transitions allow for the stimulated emission of photons, which is fundamental to the laser process.
\end{enumerate}

\section*{Long Form Response}
\begin{enumerate}[label=\arabic*.]
\setcounter{enumi}{10}
\item \textbf{Answer:} When unpolarized light passes through two ideal linear polarizers oriented at a 45-degree angle to each other, the intensity of light transmitted through the second polarizer can be calculated using Malus's Law. Initially, the intensity of unpolarized light is halved upon passing through the first polarizer, yielding an intensity of I₁ = 0.5 I₀, where I₀ is the incident intensity. The light that exits the first polarizer is now polarized, and when it encounters the second polarizer, the transmitted intensity is given by I₂ = I₁ * $cos^2($θ), where θ is the angle between the light's polarization direction and the axis of the second polarizer. With θ set at 45 degrees, the transmitted intensity becomes I₂ = (0.5 I₀) * $cos^2(45$^\ensuremath{\circ}) = (0.5 I₀) * (1/2) = 0.25 I₀. Therefore, the transmitted intensity is 25\% of the incident intensity, illustrating the principles of polarization and the effect of multiple polarizers on light intensity.
\\ \textit{Note:} correct because it accurately applies Malus's Law to determine that the first polarizer halves the intensity of unpolarized light, and the second polarizer further reduces the transmitted intensity based on the cosine of the angle between the two polarizers, leading to a final transmitted intensity of 25\% of the original incident intensity.
\item \textbf{Answer:} When a uniform solid disk rolls down an inclined plane without slipping, its total kinetic energy (KE) is distributed between translational kinetic energy and rotational kinetic energy. The translational kinetic energy, associated with the center of mass motion, is given by \( \frac{1}{2} mv^2 \), while the rotational kinetic energy, due to its spinning about its axis, is given by \( \frac{1}{2} I \omega^2 \). For a solid disk, the moment of inertia \( I \) is \( \frac{1}{2} m r^2 \), and the relationship between linear speed \( v \) and angular speed \( \omega \) is \( \omega = \frac{v}{r} \). Substituting these into the rotational energy equation reveals that the fraction of the total kinetic energy that is rotational is \( \frac{1}{3} \). Thus, one-third of the total kinetic energy of the disk is rotational, emphasizing the important interplay between translational motion and rotation in rolling objects.
\\ \textit{Note:} correct because it accurately describes how the total kinetic energy of a solid disk is divided into translational and rotational components, using the appropriate equations and the moment of inertia specific to a solid disk to illustrate the fraction of energy that is rotational.
\end{enumerate}

\vfill
\noindent\textit{End of Answer Key}
\end{document}